\scanpage{21}
\mspage{19}
logène \inserted{$C$} de $S$, un complexe de Modules ordinaire (envisagé comme objet de la catégorie dérivée des faisceaux de Modules sur $S$, resp. $T$). Un complexe de cristaux est dit pseudo-cohérent (resp. parfait, resp. $\ldots$) si pour tout objet $(U\to T)$ de $C$, le complexe induit sur $T$ est pseudo-cohérent (resp. $\ldots$). Ceci posé, voilà la théorie qu'on voudrait: A tout morphisme lisse et propre $f:X\to S$, serait associé un complexe de cristaux \inserted{(absolu)} $\operatorname{DR}(f)$ \inserted{sur $S$}, appelé cohomologie de De Rham \inserted{cristalline} de $f$. \margin{\unsure{c}} Ce complexe doit être parfait, sa formation doit être compatible avec tout changement de base sur $S$ (l'image inverse des complexes de cristaux étant entendue, bien entendu, au sens de catégories dérivées $\ldots$), et bien sûr $\operatorname{DR}(f)$ dépend fonctoriellement (de façon cotravariante) de $X$ sur $S$. Tôt ou tard, il faudra expliciter aussi une formule de Künneth $\operatorname{DR}(f\times_Sg)\simeq\operatorname{DR}(f)\overset{L}{\otimes}\operatorname{DR}(g)$, et une formule de dualité, qui pour $f$ partout de dimension relative $d$ s'exprime comme un accouplement, définissant une autodualité, $\operatorname{DR}(f)\times\operatorname{DR}(f)\longrightarrow T^d[2d]$, où $T^d$ est le cristal de Tate de poids $2d$, et où $[2d]$ indique qu'on translate les degrés de $2d$ (attention au facteur 2 !). Enfin, on veut comme de juste un isomorphisme $\operatorname{DR}(f)(S)\simeq\underline{R}f_*(\underline{\Omega}^{\bullet}_{X/S})$, fonctoriel en $X$, compatible avec les changements de base, avec Künneth et la dualité (déjà connus pour la cohomologie de De Rham ordinaire).

Par prudence, je me suis abstenu de dire quoi que ce soit sur le cas $f$ non propre, dont il faudrait parler tout au moins si on voulait faire sérieusement le lien avec Washnitzer-Monsky.

2.3. \uline{La cohomologie de De Rham comme bi\inserted{iso}cristal}. A supposer qu'on ait une théorie du type envisagé dans 2.2., \deleted{on constate aussitôt que chaque} on trouve pour chaque entier $p$ un homomorphisme de frobenius
\[
F_p:\operatorname{DR}(f)_p^{(p)}\longrightarrow\operatorname{DR}(f)_p,
\]
où l'indice $p$ au complexe de cristaux $\operatorname{DR}(f)$ désigne la restriction à $S_p=V(p.1)$ (au sens des catégories dérivées), qui d'après les conditions de 2.2. n'est autre que $\operatorname{DR}(f_p)$, $f_p:X_p\longrightarrow S_p$ étant induit par $f$. Utilisant toujours la même condi-
\margin{20}
